% =========================================================================
% Chapter 2 — Conceptual Study
% =========================================================================
\chapter{Conceptual Study}

\chapintro
Conception is the most critical phase when it comes to software development: an overall
design of the application is specified and adapted before any production-grade code is
written. For a project that combines a typed full-stack web application, three GPU AI
pipelines and a relational database with strict authorisation, getting the conception right
saves an order of magnitude of refactoring later. In this chapter we start by specifying the
design methodology, then present the global and detailed use case diagrams, the class diagram,
three sequence diagrams covering the core flows, the dataset and AI pipeline specification,
and finally the system architecture.

% -------------------------------------------------------------------------
\section{Design methodology}
\label{sec:design_methodology}

To design our system we use \textbf{UML} (Unified Modeling Language) as an object-oriented
modelling language. UML is a general-purpose, object-oriented modelling language that uses
various types of diagrams to provide a standard graphical language to visualise the design of
a system by specifying its structure, behaviour and interactions. UML helps to simplify
complex software design by reducing many words of explanation to a few graphical, more
understandable diagrams~\cite{uml_classdiag,uml_seqdiag}.

The UML diagrams used in this report are: use case diagrams (functional view of what each
actor can do), class diagram (static structure of the persistent domain), and sequence
diagrams (dynamic interactions between actors and system components for the most important
scenarios). Diagrams were drafted on draw.io and then redrawn in TikZ for inclusion in this
report.

% -------------------------------------------------------------------------
\section{Use case diagrams}
\label{sec:usecases}

Use case diagrams provide a higher-level view of the system by showing a graphical
representation of a user's various interactions. Users are represented as actors interacting
with use cases represented by ellipses. Two stereotypes are used: \texttt{<<include>>} for a
use case that is mandatorily executed as part of another (e.g. authentication is included in
every other use case), and \texttt{<<extend>>} for an optional extension.

\subsection{Use case diagram of the Content Creator}
\label{sec:uc_creator}

\subsubsection{Global use case diagram}
\label{sec:uc_creator_global}

The figure below presents the global use case diagram for the Content Creator. All use cases
include the \texttt{Authenticate} use case, mirroring the JWT-bearer-token requirement of
every protected route on the Fastify API.

\begin{figure}[H]
\centering
\begin{tikzpicture}[
  node distance = 0.5cm,
  uc/.style={ellipse,draw,minimum width=3.6cm,minimum height=0.9cm,align=center,font=\footnotesize},
  actor/.style={align=center,font=\footnotesize\bfseries},
  inc/.style={dashed,->,>=Stealth,font=\tiny},
  every edge/.style={draw}
]
  % Actor on the left (stick figure)
  \node[actor] (creator) at (0,0) {%
    \tikz{
      \draw (0,0) circle (0.18);
      \draw (0,-0.18) -- (0,-0.85);
      \draw (-0.35,-0.45) -- (0.35,-0.45);
      \draw (0,-0.85) -- (-0.25,-1.25);
      \draw (0,-0.85) -- (0.25,-1.25);
    }\\Content\\Creator};

  % System boundary
  \node[draw,rounded corners,fit={(2,2.6) (12,-3.2)},inner sep=0.5cm,label={[font=\footnotesize\itshape]above:BViral Control Center}] (sys) {};

  % Use cases
  \node[uc] (auth)      at (4.5, 2.0) {Authenticate};
  \node[uc] (upload)    at (4.5, 1.0) {Upload Video};
  \node[uc] (predict)   at (4.5, 0.0) {Get AI Score (Pillar I)};
  \node[uc] (caption)   at (4.5,-1.0) {Generate Captions (Pillar II)};
  \node[uc] (enhance)   at (4.5,-2.0) {Enhance Video (Pillar III)};

  \node[uc] (connect)   at (10.0, 2.0) {Connect Social Account};
  \node[uc] (schedule)  at (10.0, 1.0) {Schedule Post};
  \node[uc] (analytics) at (10.0, 0.0) {View Analytics};
  \node[uc] (download)  at (10.0,-1.0) {Download Result};
  \node[uc] (manage)    at (10.0,-2.0) {Manage Library};

  % Actor associations
  \draw (creator.east) -- (upload.west);
  \draw (creator.east) -- (predict.west);
  \draw (creator.east) -- (caption.west);
  \draw (creator.east) -- (enhance.west);
  \draw (creator.east) -- (download.west);
  \draw (creator.east) -- (analytics.west);
  \draw (creator.east) -- (connect.west);
  \draw (creator.east) -- (schedule.west);
  \draw (creator.east) -- (manage.west);

  % Include relationships
  \draw[inc] (upload)    -- node[above,sloped]{<<include>>} (auth);
  \draw[inc] (predict)   -- node[above,sloped]{<<include>>} (auth);
  \draw[inc] (caption)   -- node[above,sloped]{<<include>>} (auth);
  \draw[inc] (enhance)   -- node[above,sloped]{<<include>>} (auth);
  \draw[inc] (connect)   -- node[above,sloped]{<<include>>} (auth);
  \draw[inc] (schedule)  -- node[below,sloped]{<<include>>} (auth);
  \draw[inc] (analytics) -- node[below,sloped]{<<include>>} (auth);
\end{tikzpicture}
\caption{Global use case diagram of the Content Creator}
\label{fig:uc_creator}
\end{figure}

\paragraph{Detailed use case ``Upload Video and Analyse".}
After authentication, the Content Creator selects a video file (MP4 or MOV) and submits it.
The dashboard streams the file to the Fastify API via a multipart upload, the API persists
the binary in Supabase Storage, inserts a row in the \texttt{videos} table with status
\texttt{uploaded}, and enqueues a \texttt{process-video} job in the BullMQ queue. The video
worker picks it up, runs the chosen pillar, updates the status to \texttt{processing} and then
to \texttt{ready}.

\begin{figure}[H]
\centering
\begin{tikzpicture}[
  uc/.style={ellipse,draw,minimum width=3.0cm,minimum height=0.8cm,align=center,font=\footnotesize},
  actor/.style={align=center,font=\footnotesize\bfseries},
  inc/.style={dashed,->,>=Stealth,font=\tiny}
]
  \node[actor] (a) at (0,0) {Content\\Creator};
  \node[draw,rounded corners,minimum width=11cm,minimum height=4cm] (s) at (6.5,0) {};
  \node[uc] (auth) at (3.5, 1.2) {Authenticate};
  \node[uc] (sel)  at (3.5, 0.2) {Select File};
  \node[uc] (val)  at (6.5, 0.2) {Validate (size, codec)};
  \node[uc] (st)   at (9.5, 0.2) {Persist to Storage};
  \node[uc] (db)   at (6.5,-1.0) {Insert in DB \& Enqueue};
  \draw (a.east) -- (sel.west);
  \draw (a.east) -- (auth.west);
  \draw[inc] (sel) -- node[above]{<<include>>} (auth);
  \draw[inc] (sel) -- node[above]{<<include>>} (val);
  \draw[inc] (val) -- node[above]{<<include>>} (st);
  \draw[inc] (st)  -- node[right]{<<include>>} (db);
\end{tikzpicture}
\caption{Detailed use case --- Upload Video and Analyse}
\label{fig:uc_upload}
\end{figure}

\paragraph{Detailed use case ``Get AI Score (Pillar I)".}
After a video is uploaded, the creator can request a virality score. The API forwards the
request to the Pillar~I Python service, which loads the saved CLIP, Librosa, Whisper and
Llama-3 features for the video, runs the XGBoost regressor and the SHAP TreeExplainer, and
returns the score and the per-feature contributions.

\begin{figure}[H]
\centering
\begin{tikzpicture}[
  uc/.style={ellipse,draw,minimum width=3.0cm,minimum height=0.8cm,align=center,font=\footnotesize},
  actor/.style={align=center,font=\footnotesize\bfseries},
  inc/.style={dashed,->,>=Stealth,font=\tiny}
]
  \node[actor] (a) at (0,0) {Content\\Creator};
  \node[draw,rounded corners,minimum width=11cm,minimum height=4cm] (s) at (6.5,0) {};
  \node[uc] (auth) at (3.5, 1.2) {Authenticate};
  \node[uc] (sel)  at (3.5, 0.2) {Select Video};
  \node[uc] (vis)  at (6.5, 1.2) {Extract Visual + Audio};
  \node[uc] (txt)  at (9.5, 1.2) {Whisper + Llama Features};
  \node[uc] (xgb)  at (6.5, 0.2) {Run XGBoost};
  \node[uc] (shap) at (9.5, 0.2) {Compute SHAP};
  \node[uc] (rep)  at (6.5,-1.0) {Render Score \& Advice};
  \draw (a.east) -- (sel.west);
  \draw (a.east) -- (auth.west);
  \draw[inc] (sel) -- (auth);
  \draw[inc] (sel) -- (vis);
  \draw[inc] (sel) -- (txt);
  \draw[inc] (vis) -- (xgb);
  \draw[inc] (txt) -- (xgb);
  \draw[inc] (xgb) -- (shap);
  \draw[inc] (xgb) -- (rep);
\end{tikzpicture}
\caption{Detailed use case --- Get AI Score (Pillar I)}
\label{fig:uc_score}
\end{figure}

\paragraph{Detailed use case ``Generate Captions (Pillar II)".}
After authentication and selection of an uploaded video, the creator picks a caption style
(\texttt{stroke}, \texttt{yellow} or \texttt{pill}) and a number of words per flash. The
captioning worker extracts 16~kHz mono PCM audio with FFmpeg, transcribes with
\texttt{faster-whisper} (\texttt{small} model, \texttt{float16} on GPU), generates an SRT,
chunks the words two per flash, converts to ASS with the safe-font-size algorithm and burns
the subtitles in.

\begin{figure}[H]
\centering
\begin{tikzpicture}[
  uc/.style={ellipse,draw,minimum width=3.0cm,minimum height=0.8cm,align=center,font=\footnotesize},
  actor/.style={align=center,font=\footnotesize\bfseries},
  inc/.style={dashed,->,>=Stealth,font=\tiny}
]
  \node[actor] (a) at (0,0) {Content\\Creator};
  \node[draw,rounded corners,minimum width=11cm,minimum height=3.6cm] (s) at (6.5,0) {};
  \node[uc] (auth) at (3.5, 1.0) {Authenticate};
  \node[uc] (sel)  at (3.5, 0.0) {Select Style};
  \node[uc] (ext)  at (6.5, 1.0) {Extract Audio (FFmpeg)};
  \node[uc] (whisp)at (6.5, 0.0) {Transcribe (Whisper)};
  \node[uc] (ass)  at (9.5, 1.0) {SRT $\to$ ASS};
  \node[uc] (burn) at (9.5, 0.0) {Burn-in (FFmpeg)};
  \draw (a.east) -- (sel.west);
  \draw (a.east) -- (auth.west);
  \draw[inc] (sel) -- (auth);
  \draw[inc] (sel) -- (ext);
  \draw[inc] (ext) -- (whisp);
  \draw[inc] (whisp) -- (ass);
  \draw[inc] (ass) -- (burn);
\end{tikzpicture}
\caption{Detailed use case --- Generate Captions (Pillar II)}
\label{fig:uc_captions}
\end{figure}

\paragraph{Detailed use case ``Enhance Video (Pillar III)".}
The creator chooses an enhancement mode (\texttt{fast}, \texttt{hd} or \texttt{ultra}) and an
output codec/resolution. The pipeline performs quality assessment, classical filtering,
super-resolution and (optional) face restoration, then re-encodes with FFmpeg.

\begin{figure}[H]
\centering
\begin{tikzpicture}[
  uc/.style={ellipse,draw,minimum width=2.8cm,minimum height=0.8cm,align=center,font=\footnotesize},
  actor/.style={align=center,font=\footnotesize\bfseries},
  inc/.style={dashed,->,>=Stealth,font=\tiny}
]
  \node[actor] (a) at (0,0) {Content\\Creator};
  \node[draw,rounded corners,minimum width=11cm,minimum height=4cm] (s) at (6.5,0) {};
  \node[uc] (auth) at (3.3, 1.2) {Authenticate};
  \node[uc] (sel)  at (3.3, 0.2) {Choose Mode};
  \node[uc] (qa)   at (6.5, 1.2) {Quality Assessment};
  \node[uc] (filt) at (6.5, 0.2) {Classical Filters};
  \node[uc] (sr)   at (9.7, 1.2) {Real-ESRGAN SR};
  \node[uc] (face) at (9.7, 0.2) {GFPGAN Face Restore};
  \node[uc] (enc)  at (6.5,-1.0) {Encode (H.264/H.265/AV1)};
  \draw (a.east) -- (sel.west);
  \draw (a.east) -- (auth.west);
  \draw[inc] (sel) -- (auth);
  \draw[inc] (sel) -- (qa);
  \draw[inc] (qa)  -- (filt);
  \draw[inc] (filt)-- (sr);
  \draw[inc] (sr)  -- (face);
  \draw[inc] (sr)  -- (enc);
\end{tikzpicture}
\caption{Detailed use case --- Enhance Video (Pillar III)}
\label{fig:uc_enhance}
\end{figure}

\subsection{Use case diagram of the Administrator}
\label{sec:uc_admin}

After authentication, the Administrator can manage users, resolve alerts, monitor BullMQ
queues and re-run failed jobs. The figure below presents the global Administrator use case
diagram.

\begin{figure}[H]
\centering
\begin{tikzpicture}[
  uc/.style={ellipse,draw,minimum width=3.6cm,minimum height=0.9cm,align=center,font=\footnotesize},
  actor/.style={align=center,font=\footnotesize\bfseries},
  inc/.style={dashed,->,>=Stealth,font=\tiny}
]
  \node[actor] (admin) at (0,0) {Administrator};
  \node[draw,rounded corners,fit={(2,1.6)(12,-2.0)},inner sep=0.5cm] (sys) {};
  \node[uc] (auth)   at (4.5, 1.0) {Authenticate};
  \node[uc] (users)  at (4.5,-0.0) {Manage Users};
  \node[uc] (alerts) at (4.5,-1.0) {Resolve Alerts};
  \node[uc] (queue)  at (10, 1.0) {Monitor Queues};
  \node[uc] (rerun)  at (10, 0.0) {Re-run Failed Jobs};
  \node[uc] (agg)    at (10,-1.0) {Aggregate Analytics};
  \draw (admin.east) -- (auth.west);
  \draw (admin.east) -- (users.west);
  \draw (admin.east) -- (alerts.west);
  \draw (admin.east) -- (queue.west);
  \draw (admin.east) -- (rerun.west);
  \draw (admin.east) -- (agg.west);
  \draw[inc] (users)  -- node[above,sloped]{<<include>>} (auth);
  \draw[inc] (alerts) -- node[below,sloped]{<<include>>} (auth);
  \draw[inc] (queue)  -- node[above,sloped]{<<include>>} (auth);
\end{tikzpicture}
\caption{Global use case diagram of the Administrator}
\label{fig:uc_admin}
\end{figure}

% -------------------------------------------------------------------------
\section{Class diagram}
\label{sec:class_diag}

The persistent domain is implemented with the Drizzle ORM on PostgreSQL. The class diagram
in Figure~\ref{fig:class} reflects the actual schema declared under
\texttt{lib/db/src/schema/}: \texttt{users}, \texttt{accounts}, \texttt{videos},
\texttt{posts}, \texttt{analytics}, \texttt{alerts}. All five user-facing tables enable
PostgreSQL Row Level Security and declare per-action policies that restrict every operation
to the authenticated user.

\begin{figure}[H]
\centering
\begin{tikzpicture}[
  every edge/.style={draw,>=Stealth},
  font=\scriptsize,
  cls/.style={
    rectangle, draw, align=left,
    inner sep=0pt, minimum width=4.6cm,
    text width=4.6cm,
  },
  hd/.style={fill=blue!10, anchor=north, align=center,
             text width=4.6cm, inner sep=3pt, font=\scriptsize\bfseries},
  bd/.style={anchor=north, align=left,
             text width=4.6cm, inner sep=3pt, font=\scriptsize},
]
  % helper: a class is a vertical stack of header + body cells
  % USER
  \node[hd] (users-h)  at (0,3.2) {User};
  \node[bd,draw] (users-b) at (users-h.south) {+ id : UUID\\+ email : varchar(255)\\+ fullName : varchar(120)\\+ role : enum(admin, member)\\+ createdAt, updatedAt : timestamp};
  \node[draw, fit=(users-h)(users-b), inner sep=0pt] (users) {};
  % ACCOUNT
  \node[hd] (accounts-h)  at (6.5,3.2) {Account};
  \node[bd,draw] (accounts-b) at (accounts-h.south) {+ id : UUID\\+ platform : enum\\+ accountName : varchar\\+ accessToken : text (AES-256)\\+ refreshToken : text (AES-256)\\+ tokenExpiry : timestamp\\+ userId : UUID};
  \node[draw, fit=(accounts-h)(accounts-b), inner sep=0pt] (accounts) {};
  % VIDEO
  \node[hd] (videos-h) at (0,-0.5) {Video};
  \node[bd,draw] (videos-b) at (videos-h.south) {+ id : UUID\\+ userId : UUID\\+ originalUrl : text\\+ processedUrl : text\\+ duration : int\\+ status : enum(uploaded\dots failed)};
  \node[draw, fit=(videos-h)(videos-b), inner sep=0pt] (videos) {};
  % POST
  \node[hd] (posts-h) at (6.5,-0.5) {Post};
  \node[bd,draw] (posts-b) at (posts-h.south) {+ id : UUID\\+ videoId, accountId : UUID\\+ platform : enum\\+ scheduledAt, postedAt : ts\\+ status : enum(scheduled\dots failed)\\+ externalPostId : varchar};
  \node[draw, fit=(posts-h)(posts-b), inner sep=0pt] (posts) {};
  % ANALYTICS
  \node[hd] (analytics-h) at (12.8,1.2) {Analytics};
  \node[bd,draw] (analytics-b) at (analytics-h.south) {+ id : UUID\\+ postId : UUID\\+ views, likes, comments, shares : int\\+ revenue : numeric(12,2)\\+ fetchedAt : timestamp};
  \node[draw, fit=(analytics-h)(analytics-b), inner sep=0pt] (analytics) {};
  % ALERT
  \node[hd] (alerts-h) at (12.8,-2.2) {Alert};
  \node[bd,draw] (alerts-b) at (alerts-h.south) {+ id : UUID\\+ type : enum(error, copyright, success)\\+ status : enum(unresolved, resolved)\\+ message : text\\+ accountId, postId : UUID};
  \node[draw, fit=(alerts-h)(alerts-b), inner sep=0pt] (alerts) {};
  % associations
  \draw (users.east|-users-h) -- node[above,font=\tiny]{1 \dots\ *} (accounts.west|-accounts-h);
  \draw (users.south) -- node[right,font=\tiny]{1 / *} (videos.north);
  \draw (videos.east|-videos-h) -- node[above,font=\tiny]{1 \dots\ *} (posts.west|-posts-h);
  \draw (accounts.south) -- node[right,font=\tiny]{1 / *} (posts.north);
  \draw (posts.east) -- node[above,sloped,font=\tiny]{1 / *} (analytics.west);
  \draw (posts.east) -- node[below,sloped,font=\tiny]{1 / *} (alerts.west);
\end{tikzpicture}
\caption{Class diagram (PostgreSQL/Drizzle persistent domain)}
\label{fig:class}
\end{figure}

% -------------------------------------------------------------------------
\section{Sequence diagrams}
\label{sec:seq}

We document three core scenarios that exercise every major component of the system: an
upload-and-analyse flow that runs Pillar~I, a generate-captions flow that runs Pillar~II
asynchronously through BullMQ, and an enhance-video flow that runs Pillar~III.

\subsection{Sequence diagram of ``Upload and Analyse Video"}
\label{sec:seq_upload}

\begin{figure}[H]
\centering
\begin{tikzpicture}[font=\footnotesize,
  node distance = 0.6cm,
  msg/.style={->,>=Stealth},
  ret/.style={->,>=Stealth,dashed},
]
  % lifelines
  \foreach \i/\name/\x in {1/Creator/0,2/Dashboard/2.6,3/{Fastify API}/5.4,4/{BullMQ}/8.2,5/{Pillar I Worker}/10.8,6/PostgreSQL/13.4} {
    \node[draw,rectangle,minimum height=0.6cm,minimum width=2.2cm] (h\i) at (\x,0) {\name};
    \draw (\x,-0.3) -- (\x,-9.5);
  }

  \draw[msg] (0,-1)   -- node[above]{select file} (2.6,-1);
  \draw[msg] (2.6,-1.5) -- node[above]{POST /v1/videos (multipart)} (5.4,-1.5);
  \draw[msg] (5.4,-2.0) -- node[above]{insert videos (status=uploaded)} (13.4,-2.0);
  \draw[msg] (5.4,-2.5) -- node[above]{enqueue analyse-video} (8.2,-2.5);
  \draw[ret] (5.4,-3.0) -- node[above]{201 Created \{id\}} (2.6,-3.0);

  \draw[msg] (8.2,-3.5) -- node[above]{process(jobId)} (10.8,-3.5);
  \draw[msg] (10.8,-4.0) -- node[above]{update status=processing} (13.4,-4.0);
  \draw[msg] (10.8,-4.5) -- node[above]{download original} (13.4,-4.5);
  \draw[msg] (10.8,-5.0) -- node[above]{CLIP+Librosa+Whisper+Llama features} (13.4,-5.0);
  \draw[msg] (10.8,-5.5) -- node[above]{XGBoost predict + SHAP} (13.4,-5.5);
  \draw[msg] (10.8,-6.0) -- node[above]{insert ai\_scores row} (13.4,-6.0);
  \draw[msg] (10.8,-6.5) -- node[above]{update status=ready} (13.4,-6.5);
  \draw[ret] (10.8,-7.0) -- node[above]{job done} (8.2,-7.0);

  \draw[msg] (2.6,-7.7) -- node[above]{poll GET /v1/videos/\{id\}} (5.4,-7.7);
  \draw[msg] (5.4,-8.2) -- node[above]{select * from videos} (13.4,-8.2);
  \draw[ret] (5.4,-8.7) -- node[above]{score, SHAP advice} (2.6,-8.7);
  \draw[ret] (2.6,-9.2) -- node[above]{render virality dashboard} (0,-9.2);
\end{tikzpicture}
\caption{Sequence diagram --- Upload and Analyse Video (Pillar I)}
\label{fig:seq_upload}
\end{figure}

\subsection{Sequence diagram of ``Generate Captions"}
\label{sec:seq_captions}

\begin{figure}[H]
\centering
\begin{tikzpicture}[font=\footnotesize,
  node distance=0.6cm,
  msg/.style={->,>=Stealth},
  ret/.style={->,>=Stealth,dashed},
]
  \foreach \i/\name/\x in {1/Creator/0,2/Dashboard/2.5,3/{Fastify API}/5.0,4/{BullMQ}/7.5,5/{Captioning Worker}/10.5,6/Storage/13.5} {
    \node[draw,rectangle,minimum height=0.6cm,minimum width=2.0cm] (h\i) at (\x,0) {\name};
    \draw (\x,-0.3) -- (\x,-9);
  }

  \draw[msg] (0,-1) -- node[above]{Generate Captions} (2.5,-1);
  \draw[msg] (2.5,-1.5) -- node[above]{POST /v1/videos/\{id\}/captions \{style, words\}} (5.0,-1.5);
  \draw[msg] (5.0,-2.0) -- node[above]{enqueue caption-job} (7.5,-2.0);
  \draw[ret] (5.0,-2.5) -- node[above]{202 Accepted} (2.5,-2.5);

  \draw[msg] (7.5,-3.0) -- node[above]{process(jobId)} (10.5,-3.0);
  \draw[msg] (10.5,-3.5) -- node[above]{download original} (13.5,-3.5);
  \draw[msg] (10.5,-4.0) -- node[above,align=center]{ffmpeg extract\_audio (16k mono PCM)} (13.5,-4.0);
  \draw[msg] (10.5,-4.5) -- node[above,align=center]{faster-whisper transcribe (word\_timestamps=True)} (13.5,-4.5);
  \draw[msg] (10.5,-5.0) -- node[above]{generate\_srt + make\_word\_level\_srt} (13.5,-5.0);
  \draw[msg] (10.5,-5.5) -- node[above]{srt\_to\_ass\_dynamic + safe\_font\_size} (13.5,-5.5);
  \draw[msg] (10.5,-6.0) -- node[above]{burn\_ass\_subtitles (libx264 CRF18)} (13.5,-6.0);
  \draw[msg] (10.5,-6.5) -- node[above]{upload processed.mp4 + .srt} (13.5,-6.5);
  \draw[ret] (10.5,-7.0) -- node[above]{job done} (7.5,-7.0);

  \draw[ret] (5.0,-7.7) -- node[above]{download URL} (2.5,-7.7);
  \draw[ret] (2.5,-8.3) -- node[above]{captioned MP4 + SRT} (0,-8.3);
\end{tikzpicture}
\caption{Sequence diagram --- Generate Captions (Pillar II)}
\label{fig:seq_captions}
\end{figure}

\subsection{Sequence diagram of ``Enhance Video Quality"}
\label{sec:seq_enhance}

\begin{figure}[H]
\centering
\begin{tikzpicture}[font=\footnotesize,
  node distance=0.6cm,
  msg/.style={->,>=Stealth},
  ret/.style={->,>=Stealth,dashed},
]
  \foreach \i/\name/\x in {1/Creator/0,2/Dashboard/2.5,3/{Fastify API}/5.0,4/{Enhancement Worker}/8.5,5/{Real-ESRGAN+GFPGAN}/12,6/PostgreSQL/14.8} {
    \node[draw,rectangle,minimum height=0.6cm,minimum width=2.0cm] (h\i) at (\x,0) {\name};
    \draw (\x,-0.3) -- (\x,-9);
  }

  \draw[msg] (0,-1) -- node[above]{Enhance (mode=hd)} (2.5,-1);
  \draw[msg] (2.5,-1.5) -- node[above]{POST /v1/videos/\{id\}/enhance \{mode, codec\}} (5.0,-1.5);
  \draw[msg] (5.0,-2.0) -- node[above]{enqueue enhance-job} (8.5,-2.0);
  \draw[ret] (5.0,-2.5) -- node[above]{202 Accepted} (2.5,-2.5);

  \draw[msg] (8.5,-3.0) -- node[above]{preprocess (ffprobe + frame extract)} (12,-3.0);
  \draw[msg] (8.5,-3.5) -- node[above]{analyse\_quality (blur/noise/face)} (12,-3.5);
  \draw[msg] (8.5,-4.0) -- node[above]{run\_filter\_stage (auto)} (12,-4.0);
  \draw[msg] (8.5,-4.5) -- node[above]{Real-ESRGAN x2 / x4} (12,-4.5);
  \draw[msg] (8.5,-5.0) -- node[above]{GFPGAN face restoration} (12,-5.0);
  \draw[msg] (8.5,-5.5) -- node[above]{post\_process (H.264 CRF18 +faststart)} (12,-5.5);
  \draw[msg] (8.5,-6.0) -- node[above]{compute PSNR/SSIM/LPIPS} (12,-6.0);
  \draw[msg] (8.5,-6.5) -- node[above]{update videos.status=ready, processedUrl} (14.8,-6.5);
  \draw[ret] (5.0,-7.5) -- node[above]{download URL + metrics} (2.5,-7.5);
  \draw[ret] (2.5,-8.0) -- node[above]{enhanced MP4 + report} (0,-8.0);
\end{tikzpicture}
\caption{Sequence diagram --- Enhance Video Quality (Pillar III)}
\label{fig:seq_enhance}
\end{figure}

% -------------------------------------------------------------------------
\section{Dataset description and AI pipeline specification}
\label{sec:dataset}

The Pre-Posting Predictive Agent is the only pillar that requires a labelled training set.
This section documents how the dataset was collected, what features were engineered, and how
the four pillars interact at runtime.

\subsection{Data collection methodology}
\label{sec:dc}

The training dataset was scraped from \textbf{TikTok}, the platform that exposes the richest
public per-video metadata. We used the unofficial TikTok Web JSON endpoints (the same calls
the public web client makes) and a rotating proxy pool. For each scraped video we collected:

\begin{itemize}
  \item Identifiers: \texttt{video\_id}, \texttt{author\_id}, \texttt{create\_time}.
  \item Engagement counters at 7 days post-publication:
        \texttt{play\_count}, \texttt{digg\_count} (likes), \texttt{share\_count},
        \texttt{comment\_count}, \texttt{collect\_count}.
  \item Content metadata: \texttt{description}, \texttt{hashtags}, \texttt{music\_id},
        \texttt{music\_duration}, \texttt{video\_duration}.
  \item The actual video file (downloaded as an unwatermarked MP4 via the
        \texttt{playAddr} URL).
\end{itemize}

After cleaning (removing private accounts, duplicate IDs, videos longer than 60~seconds, and
videos with fewer than 200 plays --- their counters being too noisy to be useful as a label)
we retained \textbf{$\approx 18{,}500$} samples spanning four content verticals (tech,
fitness, fashion, comedy). The label is the engagement-rate-adjusted virality score:

\begin{equation}
  y = \log_{10}\!\left(1 + \frac{\text{plays} + 5\cdot(\text{likes} + \text{shares} + \text{comments})}{\text{follower\_count}+1}\right)
  \label{eq:label}
\end{equation}

The dataset was split chronologically into 70\,\% / 15\,\% / 15\,\% train / validation / test
to avoid look-ahead leakage --- a video uploaded in week $w$ may not appear in any split that
includes data from before week $w$.

\subsection{Feature engineering}
\label{sec:fe}

For each retained video the following multimodal feature blocks are extracted:

\begin{itemize}
  \item \textbf{Visual} ($D_v = 512$): one CLIP \texttt{clip-vit-base-patch32} embedding is
        computed per uniformly sampled keyframe (8 frames per clip), then averaged.
  \item \textbf{Audio} ($D_a = 7$): RMS energy, BPM, silence ratio, spectral centroid,
        spectral roll-off, zero-crossing rate and a pitch-confidence aggregate computed with
        \textsc{Librosa}.
  \item \textbf{Transcript} ($D_t = 384$): Whisper \texttt{small} transcribes the audio,
        then \texttt{all-MiniLM-L6-v2} produces a 384-dimensional sentence embedding.
  \item \textbf{Semantic LLM} ($D_l = 24$): Llama~3 8B (via Groq) is prompted with the
        transcript and produces 24 numeric features (tone polarity, emotion intensity, hook
        quality 0--1, call-to-action presence, narrative completeness, etc.).
\end{itemize}

The four blocks are concatenated to a 927-dimensional vector $\mathbf{x}$ which is then
reduced to a 473-dimensional vector $\mathbf{z}$ by \textbf{PCA}, retaining 99~\% of the
variance:

\begin{equation}
  \mathbf{z} = W^{\!\top}(\mathbf{x} - \boldsymbol{\mu}),
  \qquad W \in \mathbb{R}^{927\times 473},
  \quad \boldsymbol{\mu} \in \mathbb{R}^{927}
  \label{eq:pca}
\end{equation}

The XGBoost regressor $\hat{f}(\mathbf{z})$ is trained on $(\mathbf{z}_i, y_i)$ pairs to
minimise the squared-error objective with L2 regularisation, and SHAP TreeExplainer produces
per-feature contributions that are mapped back to human-readable feature names through the
PCA loadings.

\subsection{System architecture diagram}
\label{sec:archdiag}

The runtime architecture is summarised in Figure~\ref{fig:arch}. The dashboard talks to the
Fastify API; the API persists in PostgreSQL via Drizzle, stores binary assets in Supabase
Storage, and pushes work onto BullMQ Redis queues. Three Python workers consume those queues
and call the model libraries (CLIP, Whisper, Real-ESRGAN, GFPGAN) directly. A separate
publishing worker calls the YouTube / TikTok / Meta REST APIs.

\begin{figure}[H]
\centering
\begin{tikzpicture}[
  node distance=0.6cm,
  box/.style={draw,rounded corners,minimum width=2.6cm,minimum height=0.8cm,align=center,font=\footnotesize},
  tier/.style={draw,dashed,rounded corners,inner sep=0.45cm,font=\footnotesize\itshape},
  arr/.style={->,>=Stealth,thick},
]
  % presentation
  \node[box,fill=blue!8]   (browser) at (0,4)    {Web browser};
  \node[box,fill=blue!8]   (dash)    at (3.5,4)  {React/Vite\\Dashboard};
  \node[box,fill=blue!8]   (editor)  at (7.0,4)  {BVIRAL Video\\Editor (Next.js)};

  % business logic
  \node[box,fill=green!8]  (api)     at (0,2)    {Fastify API\\(TypeScript)};
  \node[box,fill=green!8]  (publish) at (3.5,2)  {Publishing\\Worker};
  \node[box,fill=green!8]  (queues)  at (7.0,2)  {BullMQ\\(Redis)};

  % AI services
  \node[box,fill=red!8]    (p1)      at (0, 0)   {Pillar I\\Predictive Agent};
  \node[box,fill=red!8]    (p2)      at (3.5,0)  {Pillar II\\Captioning};
  \node[box,fill=red!8]    (p3)      at (7.0,0)  {Pillar III\\Enhancement};

  % data
  \node[box,fill=yellow!15](pg)      at (0,-2)   {PostgreSQL\\(Supabase)};
  \node[box,fill=yellow!15](sb)      at (3.5,-2) {Supabase\\Storage};
  \node[box,fill=yellow!15](redis)   at (7.0,-2) {Redis};

  % external
  \node[box,fill=gray!10]  (yt)      at (10.5,3.5)  {YouTube API};
  \node[box,fill=gray!10]  (tt)      at (10.5,2.0)  {TikTok API};
  \node[box,fill=gray!10]  (meta)    at (10.5,0.5)  {Meta API};
  \node[box,fill=gray!10]  (groq)    at (10.5,-1.0) {Groq (Llama 3)};

  \draw[arr] (browser) -- (dash);
  \draw[arr] (dash) -- (editor);
  \draw[arr] (dash) -- (api);
  \draw[arr] (api)  -- (queues);
  \draw[arr] (queues) -- (p1);
  \draw[arr] (queues) -- (p2);
  \draw[arr] (queues) -- (p3);
  \draw[arr] (queues) -- (publish);
  \draw[arr] (api)  -- (pg);
  \draw[arr] (api)  -- (sb);
  \draw[arr] (queues) -- (redis);
  \draw[arr] (publish) -- (yt);
  \draw[arr] (publish) -- (tt);
  \draw[arr] (publish) -- (meta);
  \draw[arr] (p1) -- (groq);

  \node[tier, fit=(browser)(dash)(editor), label={[anchor=west]left:Presentation Tier}] {};
  \node[tier, fit=(api)(publish)(queues),  label={[anchor=west]left:Business Tier}]     {};
  \node[tier, fit=(p1)(p2)(p3),            label={[anchor=west]left:AI Services Tier}]  {};
  \node[tier, fit=(pg)(sb)(redis),         label={[anchor=west]left:Data Tier}]         {};
\end{tikzpicture}
\caption{System architecture --- 4-tier microservice layout}
\label{fig:arch}
\end{figure}

% -------------------------------------------------------------------------
\chapconclusion
In this chapter we presented the design methodology, derived the use case diagrams of the two
actors with detailed sub-diagrams for the four most important features, drew the class
diagram of the Drizzle/PostgreSQL persistent domain, walked through three sequence diagrams
covering the upload-and-analyse, captioning and enhancement flows, documented the dataset
collection and feature engineering of Pillar~I, and concluded with the four-tier system
architecture diagram. The next chapter dives into the realisation: the technology stack,
each of the four AI pillars with concrete code, evaluation metrics and screenshots.

\cleardoublepage
